% Chapter IV -- Methodology (rev L-104)
% Drafted directly from implementation evidence: key subsections name the
% scripts/configs that realize them so the panel can verify rigor that would
% otherwise remain invisible. Tense: past/present-perfect (study executed);
% proposal-stage future-tense phrasing is retired per the alignment matrix
% (internal working document; binding rules mirrored in the migration plan).
% REV L-104: dangling refs removed (fig:pipeline, tab:modelparams,
% tab:alignment); repro-map pointer corrected to Section 4.7; repo-internal
% jargon ("learned-rules file", "wiki/log.md") replaced with panel-facing
% wording; "registered" hypotheses downgraded to "predeclared in writing".

\chapter{Methodology}

\section{Chapter Overview and Design Rationale}
This chapter narrates the executed study. The investigation is a
quantitative, multi-tier benchmark study with an internal replication
structure: a \emph{pilot} established pipeline correctness, effect-size
order-of-magnitude expectations, and instrumentation stability; the
\emph{full study} re-ran the locked protocol at scale. Design choices are
reported with their alternatives and rejection reasons, because several of
them (synthetic training data, benchmark subsampling, a frozen feasibility
gate, salvaged data shards) materially shape the validity argument
assembled in Chapter~6.

Three data tiers support the study: (i) a \emph{surrogate tier} of
controlled synthetic corpora for training, (ii) an \emph{external tier}
consisting of a stratified RAID subsample for out-of-distribution
validation, and (iii) a \emph{first-party tier} of consent-gated human-
written text (including a separately quarantined tier of recovered device
shards). Each stage from data construction through training, evaluation,
and auditing is mapped to its implementing artifact in the reproducibility
table (Table~\ref{tab:repro}, Section~4.7).

\section{Research Questions and Hypotheses}
The executed study operationalizes the research questions as revised during
the proposal-to-results alignment pass (canonical wording in Chapter~1).
In condensed form:
RQ1: Can a detector trained solely on surrogate corpora achieve usable
detection (AUROC/AUPRC) on human-collected and benchmark text?
RQ2: How does surrogate-trained detection compare, under paired deltas, to
zero-shot and embedding-based baselines across generator families and
domains?
RQ3: How do calibration/reliability properties behave across strata, and
where does confidence mislead?
RQ4: What fraction of residual error is attributable to paraphrastic
evasion, length effects (headroom), and mixed authorship?
Correspondingly, directional hypotheses H1--H4 were predeclared in writing
before full-study outcome inspection; the pilot was used for
instrumentation only and is excluded from hypothesis tests.

\section{Data Construction}
\subsection{Tier S: Surrogate Synthetic Corpora (surrogate\_v2)}
Training data were manufactured rather than scraped. The surrogate\_v2
pipeline generates documents across a factorial grid spanning generator
family, prompt template, topic band, and target-length bucket, recording a
provenance tuple for every document. Rationale: (a) licensing and ethics---
scraping model outputs at scale from live services raises terms-of-service
conflicts the study avoids by construction; (b) control---strata can be
balanced deliberately rather than opportunistically; (c) reproducibility---
generation is deterministic given seed and config. Costs are equally
explicit and deferred to Section~6.2: surrogate training inherits whatever
generator idiosyncrasies the chosen families exhibit, creating a
surrogate--real gap that the external tier is designed to expose.
Generation configurations are reproduced verbatim (keys redacted) in
Appendix~C; the licensing gate that blocked non-permissive sources ran
before any corpus assembly and its verdict log is retained.

\subsection{Tier E: External Validation --- Stratified RAID Subsample}
Full RAID evaluation was infeasible under the compute budget; the subsample
was drawn by stratifying on RAID's domain, decoding strategy, and attack
condition axes with proportional allocation within strata and a floor of
$n_{\min}$ per cell [[FILL: raid subsample manifest; report exact n and
allocation]]. Subsampling is a deliberate trade: variance increases, but
coverage of the adversarial conditions central to RQ4 is preserved. The
manifest enabling exact reconstitution of the subsample ships as
Appendix~D.

\subsection{Tier F1: Own Corpus --- Consent-Gated Collection}
First-party HWT was collected prospectively. Intake forms captured
demographic and authorship metadata; the consent pipeline attached
versioned, timestamped consent receipts; the participant-data guard
screened submissions for personally identifiable content and scope
violations before admission. Splits (train/validation/test) were
partitioned at the \emph{author} level, never at the document level, to
prevent leakage of individual stylistic signatures across splits; the
partitioner's seed and assignment tables are retained. Class balance,
metadata completeness, and guard-pass rates fed the feasibility gate.

\subsection{The Frozen Feasibility Gate}
Before any model consumed Tier F1, a predeclared gate evaluated minimum
adequacy: corpus volume, HWT/MGT balance within tolerance, guard-pass
rate, and metadata completeness thresholds [[FILL: exact thresholds from
gate spec; cite freeze memo date]]. Freezing before outcome inspection is
the study's principal protection against data-dependent analytic
flexibility (the ``garden of forking paths'' concern) on the first-party
tier. The gate's verdict, its inputs, and the freeze date are reproduced
in Appendix~E. Where the gate constrained science (e.g., limiting subgroup
analyses to adequately powered strata), the constraint is stated alongside
the affected result rather than discovered post hoc.

\subsection{Tier F2: Recovered Device Shards (phd\_rescued)}
A salvage operation recovered partial writing shards from participant
devices that had failed routine upload. These enter the study as a
quarantined tier: (i) a documented chain-of-custody record per shard
(Appendix~F); (ii) re-application of the data guard; (iii) exclusion from
headline numbers unless robustness analyses show tier inclusion does not
flip any primary conclusion; (iv) dual reporting (with/without F2) for
every headline table. The provenance defense rests on demonstrating that
no analytic conclusion depends on the salvage pathway (Section~6.3).

\section{Models and Baselines}
The principal system is a discriminative detector trained on Tier S
(feature representation and architecture per run configuration
[[FILL: model config JSON; transcribe into a parameter table here]])
with early stopping on surrogate validation folds. Baselines instantiate
three families chosen to span resource requirements:
(a) \emph{Zero-shot statistical}: likelihood-ratio and perplexity-based
scorers requiring no task-specific training;
(b) \emph{Embedding-based}: frozen sentence/text encoders with lightweight
heads or distance scores;
(c) \emph{Commercial/API detectors}: queried under fixed seeds and logged
responses where terms permit, else omitted with reason recorded.
Feature-side analyses reuse the idea-phase constructs: intrinsic-dimension
estimates and linguistic-splitting slices are computed as descriptive
lenses, never as unlabeled classifiers.

\section{Evaluation Protocol}
All evaluation is item-paired. Primary metrics: AUROC (headline), AUPRC
(prevalence-robust companion), computed per stratum and pooled, with
bootstrap (item-level resampling, $B$ resamples [[FILL: metrics.py
default]]) 95\% confidence intervals. Calibration is assessed with
reliability diagrams and binned ECE per stratum. System comparisons are
made as \emph{paired deltas} on identical items with paired-bootstrap
interval estimation; multiplicity is handled by reporting all prespecified
comparisons and applying Holm correction to the family of headline tests.
Every figure in Chapter~5 regenerates from a stored run directory; the
mapping table of result artifact $\rightarrow$ generating script $\rightarrow$
configuration hash appears in Section~4.7.

\section{Robustness and Evasion Evaluation}
RQ4 is served by three audits. (1) \emph{Paraphrase evasion}: MGT test
items are passed through the evader suite (paraphrase strength levels
[[FILL: evaders config]]) before scoring; degradation curves are reported
per generator family. (2) \emph{Adversarial RAID strata}: RAID's native
attack conditions are scored as separate strata rather than pooled.
(3) \emph{Repair attempts}: where headroom analysis localizes failures
(e.g., short texts), targeted mitigations are attempted under the same
protocol and reported regardless of outcome, preventing silent
file-drawer effects.

\section{Compute Placement and Reproducibility}
All GPU-bearing steps (generation, training, large-scale scoring) executed
in cloud notebook sessions with pinned dependency sets; orchestration,
guards, and audits ran locally. Reproducibility is anchored by three
instruments: (i) immutable run directories containing config JSON, seeds,
and environment captures; (ii) a script-to-section map
(Table~\ref{tab:repro}) connecting every methods subsection to the tool
that implements it; and (iii) an append-only decision log reconstructing
the choice points, including the freeze of the feasibility gate and the
adoption of surrogate\_v2 over its predecessor.

\begin{table}[t]
\centering\small
\caption{Reproducibility map: methods subsections to implementing
artifacts (paths relative to repository root).}
\label{tab:repro}
\begin{tabular}{p{5.2cm}p{6.8cm}}
\toprule
Methods element & Implementing artifact(s) \\
\midrule
Surrogate generation grid & surrogate\_v2 pipeline scripts; generation
configs (Appendix~C) \\
Licensing gate & gate runner + verdict log \\
RAID subsampling & subsample manifest (Appendix~D) \\
Consent / guard & consent pipeline modules; guard policy (Appendices~A--B) \\
Splits and partitions & partitioner script; assignment tables \\
Feasibility gate & gate spec + freeze memo (Appendix~E) \\
Metrics (AUROC/AUPRC/ECE) & classify.py, metrics.py, embeddings.py \\
Topological slices & topology.py \\
Headroom/breakdown audits & validity\_audit\_2026-07-25 artifacts \\
Results tables/figures & pilot\_results.md, full\_study\_results.md runs \\
\bottomrule
\end{tabular}
\end{table}

\section{Ethical Considerations}
Human participation proceeded under informed consent with withdrawal rights;
collection minimized demographic capture to what analysis required. The
guard pipeline enforced privacy scrubbing pre-admission, and raw device
archives never entered version control---only guarded derivatives and
manifests, consistent with the participant-data-guard intent. Because the
thesis's subject matter includes detector-evasion techniques, the study
distinguishes sharply between (a) evasion as an \emph{object of
measurement} on experiment data and (b) any application of evasion tooling
to the researcher's own writing; the latter is prohibited outright by the
integrity firewall policy (Appendix~H), which also governs disclosure of
composition provenance for this manuscript itself.